\section{Before scaling the experiment, name its resource}
The preceding chapter followed molecular candidates through selective
operations. Before making that experiment larger, pause at a simpler question:
what exactly became faster when candidates were processed together?

Consider four candidates and three dependent predicate tests. In an abstract
machine with enough parallel test slots, all four candidates can undergo the
first test together, then the second, then the third. The elapsed number of
dependent stages is three, but twelve candidate-tests have still occurred.
The illustration counts every test, including tests on candidates that might
already have failed; early stopping would be a different work policy.

\NativeFigure{work-depth}{Rows are candidates; columns are dependent test
stages. Horizontal arrows mean precedence for the same candidate. A column
can execute in parallel only under the declared resource model. Twelve tests
and depth three are different counts, neither a measured reaction time.}{fig:textbook-work-depth}

Let $C$ be candidate count and $S$ the number of sequential stages. Under this
idealized all-candidates/all-stages policy,
\begin{equation}
W=CS\quad[\text{candidate-tests}],\qquad
D=S\quad[\text{dependent stages}].
\end{equation}
The dimensions expose the mistake in equating work with time. Converting
$D$ into seconds requires a duration model for a stage. Implementing the
parallel columns requires capacity for the candidates and the operations.
Neither requirement disappears when a diagram draws the column vertically.

\textbf{Prediction problem.} Suppose candidate count doubles while stage count
stays fixed. What happens to these two counts? Does that prove the physical
experiment takes the same time?

\textbf{Worked answer.} Work doubles from 12 to 24 candidate-tests and abstract
depth remains three. Equal physical time does not follow: concentration,
mixing, reaction completion, separation, transfers and readout may change.
The formula is a declared abstract model. The existing companion function
\texttt{work\_depth} computes its counts; it does not simulate chemistry.

There is a second, independent scaling question. A population can have many
copies without representing every relevant candidate. If one independently
drawn copy produces a detectable witness with probability $q$, then $m$ copies
miss the witness with probability $(1-q)^m$. Thus the probability of observing
at least one witness is
\begin{equation}
P_{\mathrm{signal}}=1-(1-q)^m.
\end{equation}
This calculation assumes independent copies and a fixed per-copy probability.
It is not a calibrated molecular yield model. Correlated synthesis failures
or concentration-dependent reactions can invalidate it.

The rest of the chapter keeps these distinctions separate: the logical
decision problem, the algorithmic work/depth model and the physical evidence
needed to justify a conclusion. A short sequence of stages alone cannot prove
that an exponentially large candidate population is cheap to construct.
